\documentclass[11pt,paper=a4,answers]{exam}
\usepackage{graphicx,lastpage}
\usepackage{upgreek}
\usepackage{censor}
\usepackage{gensymb}
\usepackage{amsmath}
\usepackage{multicol}
\usepackage{enumitem}
\usepackage{tasks}
\usepackage{adjustbox}
\usepackage{xcolor}
\usepackage{tfrupee}
\setlength{\voffset}{0.25in}
\usepackage{tabularx}
\newcolumntype{Y}{>{\centering\arraybackslash}X}
%==============================================================
% Customise Non-Essential Packages Below
% =============================================================
\usepackage[most]{tcolorbox}
\usepackage{eso-pic}
\usepackage{tikz}
\AddToShipoutPictureBG{%
\begin{tikzpicture}[remember picture,overlay]
\foreach \x in {-8,-4,0,4,8,12,16,20}
\foreach \y in {-8,-4,0,4,8,12,16,20} {
\node[opacity=0.05,rotate=45,anchor=center] at ([xshift=\x cm,yshift=\y cm]current page.center) {%
\begin{minipage}{2cm}
\centering
\includegraphics[width=2cm]{images/logo_black.png}\\
\small preppeo.com
\end{minipage}%
};
}
\end{tikzpicture}%
}
\printanswers

%==============================================================
% Customise Non-Essential Packages Above
% =============================================================
\definecolor{svgray}{rgb}{0.90, 0.90, 0.90}
\graphicspath{{images/}}

\usepackage[normalem]{ulem}
\renewcommand\ULthickness{1pt}   %%---> For changing thickness of underline
\setlength\ULdepth{3.5ex}%\maxdimen ---> For changing depth of underline
\renewcommand{\baselinestretch}{1}
\chead{
\includegraphics[scale=0.1,height=2cm ]{logo.png}
}
\pagestyle{empty}

\pagestyle{headandfoot}
\headrule
\cfoot{W: https://preppeo.com; M: +91-8130711689}

\begin{document}
%==============================================================
% Customise Question paper details below this
% =============================================================
\begin{center}
    \centering
    \renewcommand{\arraystretch}{1.5}
    \begin{tabularx}{\textwidth}{YYY}
    & \normalsize\bf{{\Large{{Quant}}}} & \\
    By Preppeo & \normalsize\bf{{geo\_sat\_easy\_mcq\_006}} & SAT\\
    \normalsize{{\textbf{{Algebra}} }} & \normalsize{{\textbf{{Area volumes, circle, triangle }} }} & \normalsize{{\textbf{{Solving Linear Equations}} }} \\
    \end{tabularx}
\end{center}
\par\noindent\rule{\textwidth}{1pt}
% =============================================================
% Customise Question paper details above this
% =============================================================

\begin{questions}
\pointsinrightmargin
\pointsdroppedatright
\marksnotpoints
\pointpoints{mark}{marks}
\pointformat{\boldmath\themarginpoints}
\bracketedpoints

% =============================================================
% Question Begin
% =============================================================

\section*{SAT Geometry \& Trigonometry - Practice Set 7 (Medium)}

% Q1: 30-60-90 Triangle Logic
\question
\begin{center}
\begin{tikzpicture}[scale=0.8]
    \draw[thick] (0,0) coordinate (P) -- (5.196,0) coordinate (R) -- (0,3) coordinate (Q) -- cycle;
    \draw (0,0) rectangle (0.3,0.3);
    \node[left] at (P) {$P$};
    \node[right] at (R) {$R$};
    \node[above] at (Q) {$Q$};
    \node[above right] at (2.6, 1.5) {$12$};
    \node[left] at (4.5, 0.3) {$30^\circ$};
\end{tikzpicture}
\end{center}
In the right triangle shown above, what is the length of $\overline{PQ}$?
\begin{tasks}(2)
    \task 6
    \task $6\sqrt{3}$
    \task 12
    \task 24
\end{tasks}

\begin{solution}
This is a $30^\circ-60^\circ-90^\circ$ special right triangle. In such a triangle, the side opposite the $30^\circ$ angle is exactly half the length of the hypotenuse.
\vspace{5pt}
\begin{center}
Hypotenuse $QR = 12$.
\vspace{5pt}
$PQ = \dfrac{1}{2}(QR) = \dfrac{1}{2}(12) = 6$.
\end{center}
\vspace{5pt}
Answer: \colorbox{yellow!100}{A}
\end{solution}

\vspace{1cm}

% Q2: Rectangle Diagonal (Pythagorean Theorem)
\question
The length of a rectangle's diagonal is $4\sqrt{13}$, and the length of the rectangle's shorter side is $8$. What is the length of the rectangle's longer side?
\begin{tasks}(2)
    \task 10
    \task 12
    \task $4\sqrt{3}$
    \task 144
\end{tasks}

\begin{solution}
Let the longer side be $x$. The sides and the diagonal form a right triangle, so we use the Pythagorean theorem: $a^2 + b^2 = c^2$.
\vspace{5pt}
\begin{center}
$8^2 + x^2 = (4\sqrt{13})^2$
\vspace{5pt}
$64 + x^2 = 16 \times 13$
\vspace{5pt}
$64 + x^2 = 208$
\vspace{5pt}
$x^2 = 208 - 64 = 144$
\vspace{5pt}
$x = \sqrt{144} = 12$.
\end{center}
\vspace{5pt}
Answer: \colorbox{yellow!100}{B}
\end{solution}

\vspace{1cm}

% Q3: 30-60-90 Combined Logic
\question
\begin{center}
\begin{tikzpicture}[scale=0.7]
    \draw[thick] (-3.46,0) coordinate (A) -- (2,0) coordinate (C) -- (0,3.46) coordinate (B) -- cycle;
    \draw[thick] (0,3.46) -- (0,0) coordinate (D);
    \draw (0,0) rectangle (0.3,0.3);
    \node[left] at (A) {$A$};
    \node[right] at (C) {$C$};
    \node[above] at (B) {$B$};
    \node[below] at (D) {$D$};
    \node[right] at (1, 1.8) {$8$};
    \node[left] at (-0.2, 2.5) {$30^\circ$};
    \node[left] at (1.9, 0.4) {$60^\circ$};
\end{tikzpicture}
\end{center}
In $\triangle ABC$ above, what is the length of $\overline{AD}$?
\begin{tasks}(2)
    \task 4
    \task $4\sqrt{3}$
    \task 8
    \task 12
\end{tasks}

\begin{solution}
First, look at $\triangle BDC$, which is a $30-60-90$ triangle. $BC = 8$ is the hypotenuse. The side $BD$ is opposite the $60^\circ$ angle, so $BD = \frac{8}{2}\sqrt{3} = 4\sqrt{3}$.
\vspace{5pt}
Next, look at $\triangle ABD$. It is also a $30-60-90$ triangle where $BD$ is opposite the $30^\circ$ angle. The side $AD$ is opposite the $60^\circ$ angle.
\vspace{5pt}
\begin{center}
$AD = BD \times \sqrt{3}$
\vspace{5pt}
$AD = (4\sqrt{3}) \times \sqrt{3} = 4 \times 3 = 12$.
\end{center}
\vspace{5pt}
Answer: \colorbox{yellow!100}{D}
\end{solution}

\vspace{1cm}

% Q4: Pythagorean Triple Logic
\question
A right triangle has legs with lengths of $15$ centimeters and $36$ centimeters. What is the length of this triangle's hypotenuse, in centimeters?
\begin{tasks}(2)
    \task 39
    \task 42
    \task 51
    \task 1,521
\end{tasks}

\begin{solution}
Use the Pythagorean theorem: $a^2 + b^2 = c^2$.
\vspace{5pt}
\begin{center}
$15^2 + 36^2 = c^2$
\vspace{5pt}
$225 + 1,296 = c^2$
\vspace{5pt}
$1,521 = c^2$
\vspace{5pt}
$c = \sqrt{1,521} = 39$.
\end{center}
\vspace{5pt}
Answer: \colorbox{yellow!100}{A}
\end{solution}

\vspace{1cm}

% Q5: Similar Triangles (Perimeters)
\question
The table gives the perimeters of similar triangles $DEF$ and $PQR$, where $\overline{DE}$ corresponds to $\overline{PQ}$. The length of $\overline{DE}$ is $14$.
\begin{center}
\begin{tabular}{|c|c|}
\hline
 & \textbf{Perimeter} \\ \hline
Triangle $DEF$ & 25 \\ \hline
Triangle $PQR$ & 225 \\ \hline
\end{tabular}
\end{center}
What is the length of $\overline{PQ}$?
\begin{tasks}(2)
    \task 14
    \task 112
    \task 126
    \task 214
\end{tasks}

\begin{solution}
For similar triangles, the ratio of the lengths of corresponding sides is equal to the ratio of their perimeters.
\vspace{5pt}
\begin{center}
$\dfrac{\text{Perimeter } PQR}{\text{Perimeter } DEF} = \dfrac{225}{25} = 9$
\vspace{5pt}
The scale factor is 9.
\vspace{5pt}
$PQ = DE \times 9 = 14 \times 9 = 126$.
\end{center}
\vspace{5pt}
Answer: \colorbox{yellow!100}{C}
\end{solution}

\vspace{1cm}

% Q6: Sphere Volume
\question
A sphere has a radius of $\frac{9}{2}$ feet. What is the volume, in cubic feet, of the sphere?
\begin{tasks}(2)
    \task $\frac{81\pi}{2}$
    \task $\frac{243\pi}{2}$
    \task $121.5\pi$
    \task $\frac{729\pi}{6}$
\end{tasks}

\begin{solution}
The formula for the volume of a sphere is $V = \frac{4}{3}\pi r^3$.
\vspace{5pt}
\begin{center}
$V = \dfrac{4}{3} \pi \left( \dfrac{9}{2} \right)^3 = \dfrac{4}{3} \pi \left( \dfrac{729}{8} \right)$
\vspace{5pt}
$V = \dfrac{1 \times 243}{1 \times 2} \pi = 121.5\pi$.
\end{center}
\vspace{5pt}
Answer: \colorbox{yellow!100}{C}
\end{solution}

\vspace{1cm}

% Q7: Cylinder Volume (Scale factor radius)
\question
A factory makes two sizes of cylindrical cans that both have a height of $20$ centimeters. The radius of Can X is $10$ centimeters, and the radius of Can Y is $20\%$ longer than the radius of Can X. What is the volume, in cubic centimeters, of Can Y?
\begin{tasks}(2)
    \task $2,000\pi$
    \task $2,400\pi$
    \task $2,880\pi$
    \task $4,000\pi$
\end{tasks}

\begin{solution}
First, find the radius of Can Y ($r_Y$):
\vspace{5pt}
\begin{center}
$r_Y = 10 + (0.20 \times 10) = 12 \text{ cm}$.
\vspace{5pt}
$V = \pi r^2 h = \pi (12)^2 (20) = \pi (144) (20) = 2,880\pi$.
\end{center}
\vspace{5pt}
Answer: \colorbox{yellow!100}{C}
\end{solution}

\vspace{1cm}

% Q8: Cylinder Volume from Diameter
\question
A cylinder has a diameter of $10$ inches and a height of $15$ inches. What is the volume, in cubic inches, of the cylinder?
\begin{tasks}(2)
    \task $150\pi$
    \task $375\pi$
    \task $750\pi$
    \task $1,500\pi$
\end{tasks}

\begin{solution}
The radius is half of the diameter: $r = 5 \text{ inches}$.
\vspace{5pt}
\begin{center}
$V = \pi r^2 h = \pi (5)^2 (15) = 375\pi$.
\end{center}
\vspace{5pt}
Answer: \colorbox{yellow!100}{B}
\end{solution}

\vspace{1cm}

% Q9: Radian to Degree Conversion
\question
An angle has a measure of $\frac{11\pi}{12}$ radians. What is the measure of the angle in degrees?
\begin{tasks}(2)
    \task 135
    \task 150
    \task 165
    \task 185
\end{tasks}

\begin{solution}
To convert radians to degrees, multiply the measure by $\frac{180}{\pi}$.
\vspace{5pt}
\begin{center}
$\text{Degrees} = \dfrac{11\pi}{12} \times \dfrac{180}{\pi} = 11 \times 15 = 165$.
\end{center}
\vspace{5pt}
Answer: \colorbox{yellow!100}{C}
\end{solution}

\vspace{1cm}

% Q10: Compound Angle Radians
\question
The measure of angle $A$ is $\frac{\pi}{4}$ radians. The measure of angle $B$ is $\frac{2\pi}{3}$ radians greater than the measure of angle $A$. What is the measure of angle $B$, in degrees?
\begin{tasks}(2)
    \task 120
    \task 150
    \task 165
    \task 195
\end{tasks}

\begin{solution}
First, find angle $B$ in radians:
\vspace{5pt}
\begin{center}
$B = \dfrac{\pi}{4} + \dfrac{2\pi}{3} = \dfrac{3\pi + 8\pi}{12} = \dfrac{11\pi}{12} \text{ rad}$.
\vspace{5pt}
$\text{Degrees} = \dfrac{11\pi}{12} \times \dfrac{180}{\pi} = 165$.
\end{center}
\vspace{5pt}
Answer: \colorbox{yellow!100}{C}
\end{solution}

\vspace{1cm}

% Q11: Radian Circle Geometry
\question
\begin{center}
\begin{tikzpicture}[scale=1.2]
    \draw[thick] (0,0) circle (1);
    \draw[thick] (0,0) -- (0.866,0.5) coordinate (A);
    \draw[thick] (0,0) -- (0.866,-0.5) coordinate (C);
    \node at (0,-0.2) {$O$};
    \node[above right] at (A) {$A$};
    \node[below right] at (C) {$C$};
    \draw (0.3,0.17) arc (30:-30:0.3);
    \node at (0.6, 0) {\small $\pi/3$};
\end{tikzpicture}
\end{center}
The circle above with center $O$ has a circumference of $72\pi$. What is the length of minor arc $AC$?
\begin{tasks}(2)
    \task 6
    \task 12
    \task $12\pi$
    \task $24\pi$
\end{tasks}

\begin{solution}
The ratio of the central angle to $2\pi$ (full circle) is equal to the ratio of the arc length to the circumference.
\vspace{5pt}
\begin{center}
$\text{Ratio} = \dfrac{\pi/3}{2\pi} = \dfrac{1}{6}$.
\vspace{5pt}
$\text{Arc length} = \dfrac{1}{6} \times (72\pi) = 12\pi$.
\end{center}
\vspace{5pt}
Answer: \colorbox{yellow!100}{C}
\end{solution}

\vspace{1cm}

% Q12: Cube Diagonal
\question
A cube has a side length $s = 6$ centimeters. What is the length, in centimeters, of the internal diagonal of the cube?
\begin{tasks}(2)
    \task $6\sqrt{2}$
    \task $6\sqrt{3}$
    \task 18
    \task 36
\end{tasks}

\begin{solution}
The internal diagonal of a cube with side length $s$ is calculated using the formula $d = s\sqrt{3}$.
\vspace{5pt}
\begin{center}
$d = 6\sqrt{3}$.
\end{center}
\vspace{5pt}
Answer: \colorbox{yellow!100}{B}
\end{solution}


\section*{SAT Geometry - Practice Set 10 (Medium MCQ)}

% Q13: Arc Length and Sector Logic
\question
A circle has a center at $(0,0)$ and a radius of $18$. An arc on this circle has a length of $6\pi$. What is the measure, in degrees, of the central angle subtended by this arc?
\begin{tasks}(2)
    \task $30^\circ$
    \task $60^\circ$
    \task $90^\circ$
    \task $120^\circ$
\end{tasks}

\begin{solution}
The relationship between arc length ($s$), radius ($r$), and the central angle in degrees ($\theta$) is given by the formula:
\vspace{5pt}
\begin{center}
$s = \left( \dfrac{\theta}{360} \right) \times 2\pi r$
\end{center}
\vspace{5pt}
Substitute the known values ($s = 6\pi$ and $r = 18$):
\vspace{5pt}
\begin{center}
$6\pi = \left( \dfrac{\theta}{360} \right) \times 2\pi(18)$
\vspace{5pt}
$6\pi = \left( \dfrac{\theta}{360} \right) \times 36\pi$
\end{center}
\vspace{5pt}
Divide both sides by $36\pi$:
\vspace{5pt}
\begin{center}
$\dfrac{6\pi}{36\pi} = \dfrac{\theta}{360} \implies \dfrac{1}{6} = \dfrac{\theta}{360}$
\vspace{5pt}
$\theta = \dfrac{360}{6} = 60^\circ$
\end{center}
\vspace{5pt}
Answer: \colorbox{yellow!100}{B}
\end{solution}

\vspace{1cm}

% Q14: Similar Solids and Volume
\question
Two right circular cylinders are similar. The height of the first cylinder is $5$ inches and its volume is $40\pi$ cubic inches. If the height of the second cylinder is $15$ inches, what is the volume, in cubic inches, of the second cylinder?
\begin{tasks}(2)
    \task $120\pi$
    \task $360\pi$
    \task $1,080\pi$
    \task $3,240\pi$
\end{tasks}

\begin{solution}
For similar solids, the ratio of their volumes is the cube of the ratio of their corresponding linear dimensions (scale factor $k$).
\vspace{5pt}
\begin{center}
Scale Factor $k = \dfrac{\text{Height}_2}{\text{Height}_1} = \dfrac{15}{5} = 3$
\vspace{5pt}
$\text{Volume Ratio} = k^3 = 3^3 = 27$
\vspace{5pt}
$\text{Volume}_2 = \text{Volume}_1 \times 27$
\vspace{5pt}
$\text{Volume}_2 = 40\pi \times 27 = 1,080\pi$
\end{center}
\vspace{5pt}
Answer: \colorbox{yellow!100}{C}
\end{solution}

\vspace{1cm}

% Q15: Triangle Inequality Theorem
\question
A triangle has two sides of lengths $7$ and $11$. Which of the following could be the perimeter of the triangle?
\begin{tasks}(2)
    \task 22
    \task 25
    \task 36
    \task 40
\end{tasks}

\begin{solution}
According to the Triangle Inequality Theorem, the length of the third side ($x$) must be between the difference and the sum of the other two sides.
\vspace{5pt}
\begin{center}
$(11 - 7) < x < (11 + 7)$
\vspace{5pt}
$4 < x < 18$
\end{center}
\vspace{5pt}
The perimeter ($P$) is $7 + 11 + x$, which is $18 + x$. We can find the range for the perimeter by adding $18$ to the inequality:
\vspace{5pt}
\begin{center}
$(18 + 4) < P < (18 + 18)$
\vspace{5pt}
$22 < P < 36$
\end{center}
\vspace{5pt}
The only option that falls strictly between $22$ and $36$ is $25$.
\vspace{5pt}
Answer: \colorbox{yellow!100}{B}
\end{solution}

\vspace{1cm}

% Q16: Equation of a Circle
\question
In the $xy$-plane, the graph of $(x - 4)^2 + (y + 1)^2 = 25$ is a circle. What is the length of the circle's diameter?
\begin{tasks}(2)
    \task 5
    \task 10
    \task 25
    \task 50
\end{tasks}

\begin{solution}
The standard form of a circle's equation is $(x - h)^2 + (y - k)^2 = r^2$, where $(h, k)$ is the center and $r$ is the radius.
\vspace{5pt}
\begin{center}
In this equation, $r^2 = 25 \implies r = \sqrt{25} = 5$.
\end{center}
\vspace{5pt}
The question asks for the **diameter**, which is twice the radius:
\vspace{5pt}
\begin{center}
$d = 2r = 2(5) = 10$
\end{center}
\vspace{5pt}
Answer: \colorbox{yellow!100}{B}
\end{solution}


% Q17: Arc Length and Central Angle
\question
A circle has a center at $(0,0)$ and a radius of $18$. An arc on this circle has a length of $6\pi$. What is the measure, in degrees, of the central angle subtended by this arc?
\begin{tasks}(2)
    \task $30^\circ$
    \task $60^\circ$
    \task $90^\circ$
    \task $120^\circ$
\end{tasks}

\begin{solution}
The relationship between arc length ($s$), radius ($r$), and the central angle in degrees ($\theta$) is:
\vspace{5pt}
\begin{center}
$s = \left( \dfrac{\theta}{360} \right) \times 2\pi r$
\end{center}
\vspace{5pt}
Substitute $s = 6\pi$ and $r = 18$:
\vspace{5pt}
\begin{center}
$6\pi = \left( \dfrac{\theta}{360} \right) \times 36\pi$
\vspace{5pt}
$1 = \dfrac{\theta}{360} \times 6$
\vspace{5pt}
$\dfrac{1}{6} = \dfrac{\theta}{360} \implies \theta = 60^\circ$
\end{center}
\vspace{5pt}
Answer: \colorbox{yellow!100}{B}
\end{solution}

\vspace{1cm}

% Q18: Equation of a Circle Diameter
\question
In the $xy$-plane, the graph of $(x - 4)^2 + (y + 1)^2 = 25$ is a circle. What is the length of the circle's diameter?
\begin{tasks}(2)
    \task 5
    \task 10
    \task 25
    \task 50
\end{tasks}

\begin{solution}
The standard form of a circle is $(x - h)^2 + (y - k)^2 = r^2$.
\vspace{5pt}
\begin{center}
From the equation, $r^2 = 25$, so $r = 5$.
\end{center}
\vspace{5pt}
The diameter is twice the radius:
\vspace{5pt}
\begin{center}
$d = 2r = 2(5) = 10$
\end{center}
\vspace{5pt}
Answer: \colorbox{yellow!100}{B}
\end{solution}

% =============================================================
% Question End
% =============================================================

\end{questions}
\begin{center}
\rule{.5\textwidth}{1pt}
\end{center}
\end{document}
